\subsection{Caso de característica cero: criterio jacobiano para los
morfismos diferencialmente lisos}
\label{IV.16.12}

\begin{env}[16.12.1]
\label{IV.16.12.1}
Diremos que un preesquema $X$ es \emph{de característica $p$} ($p$ igual a $0$
o a un número primo) si, para todo abierto afín $U$ de $X$, el anillo
$\Gamma(U,\mathcal{O}_X)$ es de característica $p$ (0, 21.1.1). Se deduce
de (0, 21.1.3) que, para que $X$ sea de característica $0$, es necesario y
suficiente que, para todo punto \emph{cerrado} $x$ de $X$, el cuerpo residual
$k(x)$ sea de característica $0$, o, equivalentemente, que $X$ pueda estar
provisto de una estructura de preesquema sobre $\mathbb{Q}$ (necesariamente
única).
\end{env}

\begin{theorem}[16.12.2]
\label{IV.16.12.2}
Sean $X$ un preesquema de característica $0$ y $f:X\to S$ un morfismo. Si
$\Omega_{X/S}^1$ es un $\mathcal{O}_X$-módulo localmente libre (no
necesariamente de tipo finito), $f$ es diferencialmente liso.
\end{theorem}

La cuestión es local sobre $X$, así que podemos suponer que existe una familia
$(z_\lambda)$ de secciones de $\mathcal{O}_X$ sobre $X$ tal que $(dz_\lambda)$
sea una base del $\mathcal{O}_X$-módulo $\Omega_{X/S}^1$. Aplicando el criterio
(16.11.2), basta observar que los operadores
\[
  D_{\mathbf{p}}=(\mathbf{p}!)^{-1}
  \prod_\lambda D_\lambda^{p_\lambda}
\]
(donde los $D_\lambda$ son las formas coordenadas correspondientes a la base
$(dz_\lambda)$) verifican las relaciones (16.11.2.1), lo cual se deduce del
hecho de que los $D_\lambda$ son derivaciones.

\begin{env}[16.12.3]
\label{IV.16.12.3}
El teorema anterior deja de ser cierto cuando se abandona la hipótesis de que
$X$ es de característica $0$. Por ejemplo, si
$S=\operatorname{Spec}(k)$, donde $k$ es un cuerpo de característica $p>0$,
$X=\operatorname{Spec}(K)$, donde $K=k(\alpha)$ con
$\alpha\notin k$ y $\alpha^p\in k$, se comprueba inmediatamente
\oldpage[IV-4]{56}
que $\Omega_{X/S}^1$ es de rango $1$ y que el morfismo $X\to S$ es
diferencialmente liso hasta el orden $p-1$ (16.11.3, (iii)), pero no hasta
el orden $p$. Sin embargo, la demostración de (16.12.2) prueba que, si
$\Omega_{X/S}^1$ es localmente libre y $n!1_{\mathcal{O}_X}$ es invertible en
$\Gamma(X,\mathcal{O}_X)$, entonces $X$ es diferencialmente liso sobre $S$
hasta el orden $n$.
\end{env}
